% Solutions of chapter 26 -- Flows and matchings.
\section*{Chapter \ref{ch:b2:flows-and-matchings} --- Flows and Matchings}

\begin{solution}{exo:b2:flows-and-matchings:1}
$f_1$ respects the capacities but not conservation at $c$: $5 + 5 = 10$ enter and $7$ leave. $f_2$ is balanced but puts
$9$ on $c \to t$, whose capacity is $7$; its \omterm{def:b2:flows-and-matchings:flow}{value} $14$ would also exceed the \omterm{def:b2:flows-and-matchings:cut}{cut} of capacity $12$. $f_3$ is a \omterm{def:b2:flows-and-matchings:flow}{flow}: at
$a$, $6 = 1 + 5$; at $b$, $6 = 6$; at $c$, $1 + 6 = 7$; at $d$, $5 = 5$; its \omterm{def:b2:flows-and-matchings:flow}{value} is $6 + 6 = 12$.

\omoutput{code/bachelor-2/26-flows-and-matchings/outputs/exo_is_flow.txt}{The three candidates checked.}
\end{solution}

\begin{solution}{exo:b2:flows-and-matchings:2}
\omoutput{code/bachelor-2/26-flows-and-matchings/outputs/exo_no_reverse.txt}{With and without the reverse residual
edges.}

Without reverse edges the search finds $s, a, c, t$ and $s, a, d, t$, then stops at $8$: every remaining path to $t$
needs $c \to t$, which is full, or $a \to d$ entered from $a$, which $s \to a$ can no longer feed. The \omterm{def:b2:flows-and-matchings:flow}{flow} $12$ needs
the path $s, b, c, a, d, t$, which sends $4$ units back against $a \to c$. The version misses the term $+ f(v, u)$ of
the \omterm{def:b2:flows-and-matchings:residual}{residual capacity}, at the end of line 21 of \cref{lst:b2:flows-and-matchings:network}.
\end{solution}

\begin{solution}{exo:b2:flows-and-matchings:3}
The shortest paths $s, x, z, t$ (bottleneck $4$, at $x \to z$) and $s, y, w, t$ (bottleneck $6$, at $y \to w$) give
$10$. The vertices reachable from $s$ afterwards are $s$, $x$ and $y$: $x \to y$ still has room, but $x \to z$ and $y \to
w$ are full. The \omterm{def:b2:flows-and-matchings:cut}{cut} $(\{s, x, y\}, \{z, w, t\})$ has capacity $4 + 6 = 10$; the edge $z \to y$ goes from $T$ to $S$ and
does not count.

\omoutput{code/bachelor-2/26-flows-and-matchings/outputs/exo_maxflow.txt}{The same \omterm{def:b2:flows-and-matchings:flow}{flow} and \omterm{def:b2:flows-and-matchings:cut}{cut} by program.}
\end{solution}

\begin{solution}{exo:b2:flows-and-matchings:4}
After $s, a, b, t$ the middle edge is full and its residual is $1$ backwards; $s, b, a, t$ then uses it backwards,
cancels the unit, and leaves the middle edge empty again. Each such pair raises the \omterm{def:b2:flows-and-matchings:flow}{flow} on each of the four outer
edges by $1$, and each path has bottleneck $1$, so it takes $2C$ augmentations to reach $2C$, as long as no outer edge is
full. Breadth-first search finds the two paths of two edges first, each of bottleneck $C$, and then $s$ has no residual
edge left: $2$ augmentations.
\end{solution}

\begin{solution}{exo:b2:flows-and-matchings:5}
\omoutput{code/bachelor-2/26-flows-and-matchings/outputs/exo_mincut.txt}{A few residual capacities of $f_3$, and the
vertices reachable from $s$.}

From $s$, only $s \to a$ has room ($2$); from $a$, the \omterm{def:b2:flows-and-matchings:flow}{flow} $1$ on $a \to c$ leaves $6$; from $c$, the $6$ units on $b
\to c$ can be cancelled, which reaches $b$. The edges $a \to d$ and $c \to t$ are full, and $b \to c$ is full forwards.
So $S = \{s, a, b, c\}$, and the \omterm{def:b2:flows-and-matchings:cut}{cut} $(S, \{d, t\})$ has capacity $5 + 7 = 12 = |f_3|$. By
\cref{cor:b2:flows-and-matchings:weak-duality}, $f_3$ is maximum. It is not the \omterm{def:b2:flows-and-matchings:flow}{flow} of
\cref{out:b2:flows-and-matchings:flow}: a maximum \omterm{def:b2:flows-and-matchings:flow}{flow} is not unique, its \omterm{def:b2:flows-and-matchings:flow}{value} is.
\end{solution}

\begin{solution}{exo:b2:flows-and-matchings:6}
\omoutput{code/bachelor-2/26-flows-and-matchings/outputs/exo_match_bug.txt}{Shared marks against fresh marks.}

Ada takes Monday, which is marked. Ben's only day is Monday: it is marked, so the search skips it, and Ben stays free,
although Ben, Monday, Ada, Tuesday is an augmenting path. Cy takes Tuesday, Dee takes Wednesday, and Eve's days, Monday
and Wednesday, are both marked. A mark means ``already tried in this search, without success''; it is valid only during
one search, because each augmentation changes the \omterm{def:b2:flows-and-matchings:matching}{matching} and opens new \omterm{def:b2:flows-and-matchings:alternating}{alternating paths}.
\end{solution}

\begin{solution}{exo:b2:flows-and-matchings:7}
In all three cases the program's \omterm{def:b2:flows-and-matchings:cut}{cut} is the three roads at one end: the Station's roads to places 9, 10 and 11, the
depot's roads to places 4, 7 and 24, and the Lighthouse Caf\'e's roads to places 8, 12 and 25. Each of those places has
exactly three roads, so three routes is also the obvious upper bound.

\omoutput{code/bachelor-2/26-flows-and-matchings/outputs/exo_vertex_disjoint.txt}{Routes with no place in common.}

In the split network, a route through a place uses its edge from entrance to exit, of capacity $1$, so routes of a
\omterm{def:b2:flows-and-matchings:flow}{flow} share no place, and Menger's theorem becomes: the largest number of routes with no place in common equals the
smallest number of places, other than the ends, whose removal separates them. From the Station to the Hospital the number falls to
$2$: every route passes the Museum or the Fire Station, which is what the minimum \omterm{def:b2:flows-and-matchings:cut}{cut} of the split network names. From
the Lighthouse Caf\'e to the Town Square the two places are the Museum and the Clinic. From the depot nothing changes.
\end{solution}

\begin{solution}{exo:b2:flows-and-matchings:8}
Sum, over the vertices $w \in S$, the \omterm{def:b2:flows-and-matchings:flow}{flow} out of $w$ less the \omterm{def:b2:flows-and-matchings:flow}{flow} into $w$. For $w = s$ the term is $|f|$; for the
others it is $0$ by conservation, as $t \notin S$. So the sum is $|f|$. Now group the sum by edges. An edge with both
ends in $S$ counts once positively, at its tail, and once negatively, at its head: it cancels. An edge with no end in $S$
does not appear. An edge from $S$ to $T$ appears once positively and an edge from $T$ to $S$ once negatively. What remains
is $f(S, T)$.
\end{solution}

\begin{solution}{exo:b2:flows-and-matchings:9}
\omterm{def:b2:flows-and-matchings:network}{Source} to each reader, capacity $1$ (one visit a week); reader to each of their slots, capacity $1$; slot to \omterm{def:b2:flows-and-matchings:network}{sink},
capacity its number of seats. An integral maximum \omterm{def:b2:flows-and-matchings:flow}{flow} places as many readers as possible.

\omoutput{code/bachelor-2/26-flows-and-matchings/outputs/exo_readers.txt}{The reading room with equal and unequal
slots.}

The readers $8, 16, 24, 32, 40$ have $k \equiv 0 \pmod 8$, so $k \bmod 8 = 3k \bmod 8 = 0$ and $(5k + 1) \bmod 8 = 1$:
five readers can only come at slots $0$ and $1$, which have $2 + 2 = 4$ seats. Hall's condition, with seats counted,
fails for this set; the total of 40 seats does not help, since the other slots are of no use to these readers.
\end{solution}

\begin{solution}{exo:b2:flows-and-matchings:10}
The network has capacity $1$ on the edges from $s$ and to $t$, and $|L|$ on the edges from $L$ to $R$. Suppose the
maximum \omterm{def:b2:flows-and-matchings:flow}{flow} is below $|L|$, and let $(S, T)$ be a minimum \omterm{def:b2:flows-and-matchings:cut}{cut}: its capacity is below $|L|$, so it contains no edge from
$L$ to $R$. Let $A = L \cap S$. A neighbour of a vertex of $A$ is in $S$, else an edge of capacity $|L|$ would cross the
\omterm{def:b2:flows-and-matchings:cut}{cut}: $N(A) \subseteq R \cap S$. The capacity counts the edges $s \to u$ for $u \in L \setminus A$ and $v \to t$ for $v \in
R \cap S$, so
\[
|L| > c(S, T) = |L| - |A| + |R \cap S| \geq |L| - |A| + |N(A)|,
\]
that is $|N(A)| < |A|$. This is what \cref{lst:b2:flows-and-matchings:hall} prints: Ada, Ben and Fay, on the \omterm{def:b2:flows-and-matchings:network}{source} side,
share Monday and Tuesday. Conversely, if the maximum \omterm{def:b2:flows-and-matchings:flow}{flow} is $|L|$, it is a \omterm{def:b2:flows-and-matchings:matching}{matching} of all of $L$, and the condition
holds.
\end{solution}

\begin{solution}{exo:b2:flows-and-matchings:11}
Let $f'$ follow $f$ by augmenting along a shortest path $p$ of $G_f$, and suppose some distance decreases. Among the
vertices $v$ with $\delta_{f'}(v) < \delta_f(v)$, choose one with $\delta_{f'}(v)$ smallest; $v \neq s$. Let $u$ be the
vertex before $v$ on a shortest path from $s$ to $v$ in $G_{f'}$: $\delta_{f'}(u) = \delta_{f'}(v) - 1$, and by the
choice of $v$, $\delta_f(u) \leq \delta_{f'}(u)$.

If $(u, v)$ is an edge of $G_f$, then $\delta_f(v) \leq \delta_f(u) + 1 \leq \delta_{f'}(u) + 1 = \delta_{f'}(v)$, a
contradiction. So $(u, v)$ is in $G_{f'}$ but not in $G_f$: the augmentation created it, which means it pushed \omterm{def:b2:flows-and-matchings:flow}{flow} from
$v$ to $u$, so $(v, u)$ is on $p$. Since $p$ is a shortest path in $G_f$, $\delta_f(u) = \delta_f(v) + 1$. Then
\[
\delta_f(v) = \delta_f(u) - 1 \leq \delta_{f'}(u) - 1 = \delta_{f'}(v) - 2,
\]
again a contradiction. No distance decreases.
\end{solution}

\begin{solution}{exo:b2:flows-and-matchings:12}
Call $C = (L \setminus Z) \cup (R \cap Z)$.

\emph{$C$ is a cover.} Let $u \in L$, $v \in R$ be joined by an edge. If $u \notin Z$, $u \in C$. If $u \in Z$ and $uv
\notin M$, an \omterm{def:b2:flows-and-matchings:alternating}{alternating path} reaching $u$ (empty if $u$ is free, else ending with $u$'s \omterm{def:b2:flows-and-matchings:matching}{matching} edge) extends by $uv$,
so $v \in Z$ and $v \in C$. If $u \in Z$ and $uv \in M$, then $u$ is matched, so it is not one of the starting vertices,
and it was reached through its \omterm{def:b2:flows-and-matchings:matching}{matching} edge, from $v$: again $v \in Z$.

\emph{$|C| = |M|$.} Every vertex of $L \setminus Z$ is matched, since the free vertices of $L$ are in $Z$. Every vertex of
$R \cap Z$ is matched, otherwise the \omterm{def:b2:flows-and-matchings:alternating}{alternating path} reaching it would be augmenting, and $M$ would not be maximum
(\cref{thm:b2:flows-and-matchings:berge}). No edge of $M$ has both ends in $C$: if $v \in R \cap Z$ is matched to $u$,
then $u \in Z$. So $v \mapsto$ its \omterm{def:b2:flows-and-matchings:matching}{matching} edge is injective on $C$, $|C| \leq |M|$, and $|C| \geq |M|$ for any cover.

\omcode{code/bachelor-2/26-flows-and-matchings/listings/exo_cover.py}{6}{23}{The cover read from a maximum \omterm{def:b2:flows-and-matchings:matching}{matching}}

\omoutput{code/bachelor-2/26-flows-and-matchings/outputs/exo_cover.txt}{The six readers: a cover of size $5$, and no
smaller one.}
\end{solution}

\begin{solution}{pb:b2:flows-and-matchings:1}
\textbf{1.} A two-way road gives two edges, a one-way road one: $2 \times 52 + 1 = 105$, so one road is one-way (road 1,
from the Town Square to the Library).

\textbf{2.} $30$, $50$ and $70$ km/h give $20$, $40$ and $60$. Yes: the depot's road to the South Bridge carries $60$,
and so do several roads on the third, fourth and fifth routes.

\textbf{3.} The two directions of a road are separate lanes. With a shared capacity, $f(u, v) + f(v, u) \leq c$, the
network would be undirected; by question 4 a maximum \omterm{def:b2:flows-and-matchings:flow}{flow} never needs both directions, so the shared and the separate
models have the same maximum here.

\textbf{4.} If $f(u, v) > 0$ and $f(v, u) > 0$, subtract their minimum from both: capacities still hold, every vertex
stays balanced and the \omterm{def:b2:flows-and-matchings:flow}{value} does not change. The listing does it already, cancelling reverse \omterm{def:b2:flows-and-matchings:flow}{flow} before adding.

\textbf{5.} Add a vertex for the car park with an edge of capacity $100$ from the ground, and make it the \omterm{def:b2:flows-and-matchings:network}{sink}. The \omterm{def:b2:flows-and-matchings:cut}{cut}
that separates it alone has capacity $100$, and the \omterm{def:b2:flows-and-matchings:flow}{flow} of $120$ can be reduced to $100$, so the maximum is $100$.

\textbf{6.} The number of buses an hour that leave the depot, which is the number arriving at the ground. At the Town
Square, as many buses an hour arrive as leave.

\textbf{7.} No road joins the depot to the ground, so the shortest routes have two roads, and breadth-first search
finds a shortest augmenting path each time; 16-4-15 and 16-7-15 are the two of length 2.

\textbf{8.} After the first two, $4 \to 15$ carries $20$ of $20$, $7 \to 15$ and $16 \to 7$ carry $40$ of $40$, and
$16 \to 4$ has $20$ left. The path enters by $16 \to 4$ and cannot end with $4 \to 15$ or $7 \to 15$, which are full.

\textbf{9.} $2, 2, 9, 10, 14$: by \cref{lem:b2:flows-and-matchings:monotone} the distance of the ground never decreases,
and each path is a shortest path.

\textbf{10.} $|V| \cdot |E| = 30 \times 105 = 3150$ against the capacity out of the depot, $40 + 40 + 60 = 140$; the
second is much better here, the first does not depend on the capacities.

\textbf{11.} Yes: any choice of augmenting paths ends with a maximum \omterm{def:b2:flows-and-matchings:flow}{flow} (\cref{thm:b2:flows-and-matchings:maxflow-mincut}),
and with integer capacities each raises the \omterm{def:b2:flows-and-matchings:flow}{value} by at least $1$, so at most $120$ augmentations.

\textbf{12.} $20$ on road 11 and $40$ on road 20, their capacities: the edges from $S$ to $T$ of the \omterm{def:b2:flows-and-matchings:cut}{cut} built from the
\omterm{def:b2:flows-and-matchings:residual}{residual network} are saturated (proof of \cref{thm:b2:flows-and-matchings:maxflow-mincut}).

\textbf{13.} $40 + 20 + 20 + 40 = 120$. A route starts in $S$ and ends in $T$, so at its first step from $S$ to $T$ it
uses an edge from $S$ to $T$, one of the four.

\textbf{14.} Road 16 lies in this minimum \omterm{def:b2:flows-and-matchings:cut}{cut} but not in the \omterm{def:b2:flows-and-matchings:cut}{cut} nearest the ground, whose capacity stays $120$; the
\omterm{def:b2:flows-and-matchings:flow}{flow} is bounded by every \omterm{def:b2:flows-and-matchings:cut}{cut}, so it cannot pass $120$. Lying in one minimum \omterm{def:b2:flows-and-matchings:cut}{cut} does not make a road a bottleneck.

\textbf{15.} It starts from the ground and adds every place from which some place already found can be reached in the
\omterm{def:b2:flows-and-matchings:residual}{residual network}; the \omterm{def:b2:flows-and-matchings:cut}{cut} puts these places on the \omterm{def:b2:flows-and-matchings:network}{sink} side. Here only the ground itself is found, so the \omterm{def:b2:flows-and-matchings:cut}{cut} is the
five roads into the ground: $20 + 40 + 20 + 20 + 20 = 120$.

\textbf{16.} If a road $e$ is missing from some minimum \omterm{def:b2:flows-and-matchings:cut}{cut} $(S, T)$, widening $e$ does not change $c(S, T)$, which still
bounds the \omterm{def:b2:flows-and-matchings:flow}{flow}: the maximum stays the same. Roads 11 and 20 are in both \omterm{def:b2:flows-and-matchings:cut}{cuts} printed, and widening them gives 140.

\textbf{17.} Yes, with integer capacities and a widening of at least $1$: every minimum \omterm{def:b2:flows-and-matchings:cut}{cut} grows by the widening, and
every other \omterm{def:b2:flows-and-matchings:cut}{cut} had capacity at least the minimum plus $1$, so the new minimum, hence the new maximum \omterm{def:b2:flows-and-matchings:flow}{flow}, is larger.

\textbf{18.} The three roads out of the depot, roads 38, 12 and 21, with $60 + 40 + 40 = 140$, the bound of question 10.
The result of the problem: max \omterm{def:b2:flows-and-matchings:flow}{flow} $=$ min \omterm{def:b2:flows-and-matchings:cut}{cut} $= 120$ buses an hour, and the roads in every minimum \omterm{def:b2:flows-and-matchings:cut}{cut}, 11 and 20,
are the ones to widen, by $20$ at most before the depot's own roads limit the \omterm{def:b2:flows-and-matchings:flow}{flow}.

\textbf{19.} Double-deck ($i \bmod 3 = 1$): $B_1, B_4, B_7$, drivers Bo, Dia, Fen, Hal. Electric ($2$): $B_2, B_5, B_8$,
drivers Cal, Eli, Hal. Single-deck ($0$): $B_3, B_6, B_9$, drivers Ana, Bo, Eli, Gus, Ivo.

\textbf{20.} Each printed pair has the right licence, and each driver and each bus appear once: a perfect \omterm{def:b2:flows-and-matchings:matching}{matching}.

\textbf{21.} The electric buses $\{B_2, B_5, B_8\}$ have only Cal and Eli as neighbours: three buses, two drivers.

\textbf{22.} Cal, Eli and the six buses $B_1, B_3, B_4, B_6, B_7, B_9$: the electric edges end at Cal or Eli, the others
at a double-deck or single-deck bus. Its size $8$ equals the \omterm{def:b2:flows-and-matchings:matching}{matching}, which proves both optimal.

\textbf{23.} A new driver with the electric licence: the \omterm{def:b2:flows-and-matchings:matching}{matching} becomes Cal, Eli and the new driver on the three
electric buses. A driver with another licence only adds edges to buses already in the cover of question 22, which stays
a cover of size 8, so the maximum stays 8: the electric licence is the only answer.

\textbf{24.} \omterm{def:b2:flows-and-matchings:network}{Source} to each driver, capacity $2$; driver to each bus they may drive, capacity $2$; bus to \omterm{def:b2:flows-and-matchings:network}{sink}, capacity
$3$. An integral maximum \omterm{def:b2:flows-and-matchings:flow}{flow} gives the number of trips of each driver on each bus.
\end{solution}
